% NORMALIZACIÓN DE TRANSCRIPCIÓN: los superíndices mecanografiados con letra o
% se representan como cero; no hay diferencias materiales entre los testigos
% francés diplomático y corregido en este componente.

\subsection*{3.7. Compatibilidades diversas}

En esta sección examinamos el comportamiento, al variar $P$ y $Q$,
de la equivalencia de categorías cofibradas de 3.6.4. La varianza de
la categoría cofibrada de las biextensiones de $(P,Q)$ por un $G$
variable, con respecto a $P,Q$ variables, se definió en principio en
2.4; también puede expresarse diciendo que se tiene un funtor
cofibrado (2.6.1). Obsérvese que la varianza en $P,Q$ es
\underline{contravariante}. La varianza de la categoría cofibrada
$\underline\Psi_{L_\bullet(P,Q)}$ con respecto a $P,Q$ variables
también está clara, puesto que
$L_\bullet(P,Q)=L_\bullet(P)\otimes L_\bullet(Q)$ depende
funtorialmente del par $(P,Q)$, mientras que
$\underline\Psi_{L_\bullet}$ depende contravariantemente de
$L_\bullet$ (3.4). Sentado esto, una primera compatibilidad se expresa,
para toda terna de morfismos
\begin{equation}\tag{3.7.1}
 P'\longrightarrow P,\qquad
 Q'\longrightarrow Q,\qquad
 G\longrightarrow G',
\end{equation}
por el hecho de que el diagrama de funtores
\begin{equation}\tag{3.7.2}
\begin{tikzcd}[column sep=huge,row sep=large]
 \underline\Psi_{L_\bullet(P,Q)}(G)
   \arrow[r]\arrow[d]
 &\underline\Psi_{L_\bullet(P',Q')}(G')
   \arrow[d]\\
 \operatorname{BIEXT}(P,Q;G)
   \arrow[r]
 &\operatorname{BIEXT}(P',Q';G')
\end{tikzcd}
\end{equation}
es conmutativo salvo un isomorfismo canónico. Admitiremos esta
compatibilidad: explicitar el isomorfismo de conmutatividad solo podría
ser tan evidente como fastidioso (al igual, por lo demás, que la
explicitación completa que presupone de la equivalencia de categorías
cofibradas de 3.6.4).

3.7.3. De la compatibilidad precedente resulta, en particular, que el
isomorfismo (3.6.5.1) es funtorial no solo en $G$, sino también en $P$
y $Q$. Análogamente, y de manera más trivial, puede enunciarse la
funtorialidad análoga del isomorfismo definido mediante (3.1.8), del
mismo modo que (3.6.5.1) se definió mediante (3.2.4):
\begin{equation}\tag{3.7.4}
 \Biext^0(P,Q;G)
 \xrightarrow{\sim}
 \Ext^0(P\mathbin{\overset{\mathbb L}{\otimes}}Q,G)
 \simeq\Hom(P\otimes Q,G).
\end{equation}
Pero ya se había obtenido un isomorfismo de este tipo en (2.5.4.1), y
en 2.6 se observó que era funtorial en sus tres argumentos. Es
inmediato comprobar que los dos isomorfismos (3.7.4) y (2.5.4.1) son
iguales.

3.7.5. El interés del isomorfismo de 3.6.5 radica en que expresa el
notable grupo de origen geométrico $\Biext^1(P,Q;G)$ en términos del
invariante muy sencillo de álgebra homológica
$\Ext^1(P\mathbin{\overset{\mathbb L}{\otimes}}Q,G)$. Este último
forma parte de la sucesión de invariantes
$\Ext^i(P\mathbin{\overset{\mathbb L}{\otimes}}Q,G)$, a la que se
aplica el formalismo habitual del álgebra homológica. En particular,
la sucesión de estos $\Ext^i$ es un \underline{funtor cohomológico} en
el sentido de [3] (o una \underline{sucesión conexa de funtores},
en la terminología de [5]) en cada uno de los tres argumentos
$P,Q,G$. Así, una sucesión exacta
\begin{equation}\tag{3.7.5.1}
 0\longrightarrow G'\longrightarrow G\longrightarrow G''
 \longrightarrow0
\end{equation}
(respectivamente,
\begin{equation}\tag{3.7.5.2}
 0\longrightarrow P'\longrightarrow P\longrightarrow P''
 \longrightarrow0,\qquad\text{o bien}\qquad
 0\longrightarrow Q'\longrightarrow Q\longrightarrow Q''
 \longrightarrow0)
\end{equation}
da lugar a una sucesión exacta infinita de los $\Ext^i$. Sus seis
primeros términos, en los que solo intervienen $\Ext^0$ y $\Ext^1$,
pueden interpretarse en términos de biextensiones mediante 3.6.5 y
3.7.4. Ya señalamos en 3.3.6 una interpretación geométrica de la
sucesión exacta deducida de (3.7.5.1), y en especial del homomorfismo
de borde
\begin{equation}\tag{3.7.5.3}
 \partial:\Biext^0(P,Q;G'')\longrightarrow
 \Biext^1(P,Q;G'),
\end{equation}
en el marco de una \underline{categoría cofibrada exacta a izquierda}.
Puede buscarse una interpretación análoga de la sucesión exacta de
seis términos asociada a cualquiera de las sucesiones exactas cortas
de (3.7.5.2), digamos la primera; o, lo que equivale a lo mismo en
vista de 3.7.3, una interpretación geométrica del homomorfismo de borde
\begin{equation}\tag{3.7.5.4}
 \partial:\Biext^0(P',Q;G)\longrightarrow
 \Biext^1(P'',Q;G).
\end{equation}
En efecto, se tiene el resultado siguiente.

\underline{Proposición} 3.7.6. \underline{La categoría cofibrada sobre}
\[
 \operatorname{GRAB}(\underline T)^{\mathrm{op}}
 \times\operatorname{GRAB}(\underline T)^{\mathrm{op}}
 \times\operatorname{GRAB}(\underline T)
\]
\underline{de las biextensiones de $(P,Q)$ por $G$, con $P,Q,G$
variables} (cf. (2.6.1)), \underline{es exacta a izquierda no solo en
el argumento $G$, sino también en cada uno de los argumentos $P$ y
$Q$. Además, el homomorfismo de borde (3.7.5.4) deducido de la
sucesión exacta de los $\Ext^i$ asociada a (3.7.5.2), mediante los
isomorfismos 3.6.5 y 3.7.4, es el homomorfismo de borde de la teoría de
las categorías cofibradas exactas a izquierda sobre}
$\underline C=\operatorname{GRAB}(\underline T)^{\mathrm{op}}$
\underline{recordado en 3.3.6, aplicado a la categoría cofibrada de
las biextensiones de $(P,Q)$ por $G$, donde $Q$ y $G$ están fijos y
$P$ varía} (salvo un posible error de signo en 3.4.9; cf. 3.4.9.2).

Esbocemos la demostración de 3.7.6. Mediante 3.6.4 y la compatibilidad
(3.7.2), interpretamos las biextensiones de $(P,Q)$ por $G$ como
objetos de $\underline\Psi_{L_\bullet(P,Q)}(G)$. Consideremos el
homomorfismo
\[
 L_\bullet(P,Q)
 =L_\bullet(P)\otimes L_\bullet(Q)
 \longrightarrow P\otimes L_\bullet(Q)
\]
definido por el aumento $L_\bullet(P)\longrightarrow P$. De 3.5.3 y
de la planitud de $L_\bullet(Q)$ se sigue que este homomorfismo induce
isomorfismos sobre $H_0$ y $H_1$ y, por tanto, mediante 3.4.7, una
equivalencia
\[
 \underline\Psi_{P\otimes L_\bullet(Q)}
 \longrightarrow
 \underline\Psi_{L_\bullet(P,Q)}.
\]
Podemos, pues, interpretar las biextensiones de $(P,Q)$ por $G$ como
objetos de $\underline\Psi_{P\otimes L_\bullet(Q)}(G)$. La sucesión
exacta (3.7.5.2) da ahora una sucesión exacta de complejos de cadenas
\begin{equation}\tag{3.7.6.1}
 0\longrightarrow P'\otimes L_\bullet(Q)
 \longrightarrow P\otimes L_\bullet(Q)
 \longrightarrow P''\otimes L_\bullet(Q)
 \longrightarrow0.
\end{equation}
La exactitud a izquierda enunciada resulta de aplicar 3.4.8 a esta
sucesión exacta de complejos. (La comprobación directa sería, en todo
caso, inmediata.) Además, el homomorfismo (3.7.5.4) es precisamente el
homomorfismo de borde $\Ext^0\longrightarrow\Ext^1$ deducido de esta
sucesión exacta al tomar los $\Ext^i$ con valores en $G$. La última
afirmación de 3.7.6 es, por consiguiente, un caso particular de 3.4.9.
